\scanpage{197}

Or à priori on a

\[ \xi^{n}=\sum_i c_{\unsure{q|p|n}i}\,e_{0\,i} \]

les $c_{qi}\in\Lambda^+$, d'où en \unsure{transformant} par $i^*$ :

\[ e_2^n=\sum_i i^*(c_{qi})\,e_i \]

\deleted{Or \unsure{on a} $i^*(c_{p}$}

Or on a

\[ i^*(\lambda\deleted{\unsure{p}})=d(\lambda)\qquad\lambda\in\Lambda \]

$d$ étant l'augmentation de $\Lambda$, d'où

\[ e_2^n=\sum_i d(c_{qi})\,e_{\unsure{i}} \]

Or on sait déjà qu'on a $e_2^q=\sum_1^{q+1}a_{qi}\,e_i$, les $a_{qi}$ étant des entiers $\geq0$ et \deleted{\unsure{tels}} $a_{q,q+1}=1$. Donc on a

\[ \varepsilon(c_{\deleted{\unsure{n|p}}qi})=0\quad\text{si}\quad i>q+1,\qquad \varepsilon(c_{q,q+1})=1 \]

d'où, en raison des $c_{qi}\in\Lambda^+$,

\[ c_{qi}=0\quad\text{si}\quad i>q+1,\qquad c_{q,q+1}\ \text{de la forme}\ \sigma_\alpha \]

($\alpha$ dépendant de $n$ \deleted{\unsure{$q$}}). D'où les formules (1), mais on doit encore déterminer \deleted{\unsure{$=$}} \deleted{\unsure{$c_{p,p+1}=\sigma_\alpha$}.} \deleted{\illegible{3}} $c_{p,p+1}=\sigma_\alpha$

Or le plus haut poids dans la repr. $\sigma_\alpha e_{0\,p+1}=e_{\alpha,p+1}$ est $\alpha+pn$, \deleted{et} le plus

% archive sheet number: 196
